
\begin{definition}
	Функция $f$, регулярная на $\Cm$, называется \textit{целой}.
\end{definition}

\begin{note}
	Каждая целая функция представима рядом Тейлора на всей плоскости $\Cm$:
	\[f(z) = \sum_{n=0}^{+\infty}c_nz^n,~z \in \Cm\; (R_{\text{сх.}} = \infty)\]
\end{note}

\begin{proposition}
	Пусть выполнены следующие условия:
	\begin{enumerate}
		\item Функция $f$ "--- целая
		\item Существуют числа $R_0 > 0$, $A > 0$ и $m \in \N \cup \{0\}$ такие, что на $\{z \in \Cm: |z| > R_0\}$ выполнено неравенство $|f(z)| \le A|z|^m$
	\end{enumerate}
	
	Тогда $f$ является многочленом степени не выше $m$.
\end{proposition}

\begin{proof}
	Для каждого $R > R_0$ положим $M(R) := \max\{|f(z)|: |z| = R\}$. По неравенству Коши для коэффициентов ряда Лорана, для каждого $n >m $ имеем:
	\[|c_n| \le \frac{M(R)}{R^n} \le \frac{AR^m}{R^n} \xrightarrow[R \to +\infty]{} 0\]
	
	Значит, $f = \sum\limits_{n = 0}^mc_nz^n$.
\end{proof}

\begin{corollary}[теорема Лиувилля]
	Если $f$ "--- целая функция, и для некоторого $R_0 > 0$ она ограничена на $\{z \in \Cm: |z| > R_0\}$, то $f$ постоянна на $\Cm$.
\end{corollary}

\begin{proof}
	Применим утверждение выше для случая, когда $m = 0$.
\end{proof}

\begin{theorem}[основная теорема алгебры]
	Любой многочлен $P \in \Cm[z]$, степень которого не меньше 1, имеет корень в $\Cm$.
\end{theorem}

\begin{proof}[Доказательство с помощью теоремы Лиувилля]
	Предположим противное, тогда функция $\frac 1P$ "--- целая, причем выполнено $\lim\limits_{z \to \infty} \frac1P(z) = 0$, то есть эта функция ограничена в $\Cm$. Тогда, по теореме Лиувилля, $\frac 1P$ постоянна на $\Cm$, но тогда и функция $P$ постоянна на $\Cm$, что неверно.
\end{proof}

\begin{corollary}
    Из этого легко следует с помощью теоремы Безу, что многочлен $n$-ой степени имеет \(n\) корней с учётом их кратности.
\end{corollary}